\documentclass[11pt]{article}
\usepackage[margin=.78in,top=.65in,bottom=.65in]{geometry}
\usepackage{amsmath,amssymb,amsthm,hyperref}
\renewcommand{\normalsize}{\fontsize{10}{12}\selectfont}
\DeclareMathSizes{10}{10}{8}{6}
\hypersetup{colorlinks=true,linkcolor=blue,urlcolor=blue,citecolor=blue}
\setlength{\emergencystretch}{2em}
\newtheorem{theorem}{Theorem}
\newtheorem{lemma}[theorem]{Lemma}
\newcommand{\R}{\mathbb R}
\begin{document}
\begin{center}
{\Large\bfseries The adjacent-count inequality for the flat Bernoulli game}\\[.5em]
Research proof and Lean formalization\\
October 27, 2026
\end{center}
\begin{abstract}
We prove the adjacent-count inequality in Hetyei's Conjecture~6.10.
After factorial normalization, the counting sequence is a renewal sequence
whose increments are negative from index 11 onward. Finite symmetric
arrowhead matrices realize approximating moments. Their spectrum has one
negative value, and positive-weight power inequalities propagate three exact
coefficient checks to the entire tail. Classical finite power sums identify
the limiting coefficients. The Lean proof includes the source's finite
ordered-composition formula and all spectral, recurrence and limit bridges.
\end{abstract}

\section{The count formula and the target}
Hetyei~\cite[(16), Conjecture~6.10]{H} defines
\begin{equation}\label{counts}
 K_n=\sum_{\substack{m_1+\cdots+m_k=n\\m_i\ge2}}
       \frac{n!}{m_1!\cdots m_k!}\prod_{i=1}^k(m_i-2)!,
 \qquad \kappa_{n+1}=(n+1)K_n.
\end{equation}
The sum ranges over all ordered compositions; the empty composition gives
$K_0=1$, consistently with the source's generating function and
Proposition~6.7. The source proves the interpretation of these formulas as
game-position counts. Our formalization starts from the explicit finite formula.
\begin{theorem}\label{main}
For every integer $n\ge11$, $K_n<nK_{n-1}$. Equivalently, for every
integer $n\ge12$, $\kappa_n<n\kappa_{n-1}$.
\end{theorem}
Put $u_n=K_n/n!$ and $b_n=u_n-u_{n-1}$ for $n\ge1$.
Splitting a composition at its first block gives
\begin{equation}\label{rec}
 u_0=1,\quad u_1=0,\qquad
 u_{n+2}=\sum_{k=0}^{n}\frac{u_{n-k}}{(k+2)(k+1)}.
\end{equation}
This also follows from the source's generating function
$\sum u_nt^n=((1-t)(1-\log(1-t)))^{-1}$.
Theorem~\ref{main} is exactly $b_n<0$ for $n\ge11$.
Exact rational arithmetic in~\eqref{rec} gives
\begin{equation}\label{base}
 b_3=-\frac13,\qquad b_{12}=-\frac{10819}{9979200}<-\frac1{2000},
 \qquad b_{14}-b_{12}=-\frac{1707491}{1816214400}<0.
\end{equation}
We propagate these finite checks by a spectral argument, rather than
extrapolating the observed signs.

\section{Finite matrix realization}
For $N\ge1$, set $t_i=i/N$, $r_i=N^{-1/2}$, $s_i=\sqrt{1-t_i}$,
$v_i=s_ir_i$ for $0\le i<N$, and form
\[
 H_N=\begin{pmatrix}0&v^{\mathsf T}\\v&\operatorname{diag}(t_i)\end{pmatrix},
 \qquad u_n^{(N)}=(H_N^n)_{00},\quad b_n^{(N)}=u_n^{(N)}-u_{n-1}^{(N)}.
\]
For $x=(a,z)$, direct expansion gives
\begin{align*}
 x^{\mathsf T}(I-H_N)x&=\sum_i(s_iz_i-r_ia)^2,\\
 x^{\mathsf T}(I+H_N)x&=\sum_i(s_iz_i+r_ia)^2+2\sum_it_iz_i^2.
\end{align*}
Thus every eigenvalue lies in $[-1,1]$. The quadratic form of $H_N$
is nonnegative on the codimension-one subspace $a=0$.

\newpage
\section{The spectral propagation argument}
There is at most one negative eigenvalue, counted with multiplicity.
Indeed, each negative eigenvector has nonzero center coordinate, and two
orthonormal negative eigenvectors would yield a nonzero linear combination
with center coordinate zero and negative quadratic form. Let $w_\lambda$
be the squared center coordinates of an orthonormal eigenbasis. Then
\[
 \sum w_\lambda=1,\qquad \sum w_\lambda\lambda=0,\qquad
 \sum w_\lambda\lambda^2=\frac{N+1}{2N}\ge\frac12.
\]
Consequently a negative eigenvalue exists. Write it as $-q$, where $0<q\le1$.
Since $(\lambda+q)(1-\lambda)\ge0$ at every spectral value,
$\lambda^2\le(1-q)\lambda+q$; averaging gives $q\ge1/2$.
The center-moment expansion yields
\begin{equation}\label{inc}
 b_n^{(N)}=\sum_\lambda w_\lambda(\lambda-1)\lambda^{n-1}.
\end{equation}
The following finite criterion is the main step.
\begin{lemma}\label{prop}
Suppose a finite positive spectral measure has mass~1, mean~0, support in
$\{-q\}\cup[0,1]$, and $1/2\le q\le1$. Write $b_n$ for~\eqref{inc}.
If $b_{12}\le-c<0$, $b_{14}\le b_{12}$ and $b_3\le-1/4$, then for every
integer $m\ge0$,
\[
 b_{12+2m}\le-c\,4^{-m},\qquad b_{2m+1}\le-4^{-m}.
\]
\end{lemma}
\begin{proof}
For the even indices put
\[
 T=w_{-q}(1+q)q^{11},\quad
 W_\lambda=w_\lambda(1-\lambda)\lambda^{11},\quad
 S_m=\sum_{\lambda\ge0}W_\lambda\bigl((\lambda/q)^2\bigr)^m.
\]
Then $b_{12+2m}=q^{2m}(T-S_m)$. The hypotheses and $q\le1$ give
$q^2(T-S_1)=b_{14}\le b_{12}\le q^2b_{12}=q^2(T-S_0)$,
hence $S_1\ge S_0$. The elementary inequality
$z^m\ge1+m(z-1)$ for $z\ge0$ implies $S_m\ge S_0$.
Thus $b_{12+2m}\le-cq^{2m}\le-c4^{-m}$.
For odd indices, $d_\lambda=w_\lambda(1-\lambda)$ are nonnegative,
$\sum d_\lambda=1$, and $\sum d_\lambda\lambda^2=-b_3\ge1/4$.
The same power inequality, normalized by $1/4$, gives
$-b_{2m+1}=\sum d_\lambda(\lambda^2)^m\ge4^{-m}$.
\end{proof}

\section{Exact coefficient limits and strictness}
Let $A_n=u_n^{(N)}$ and let $z_{n,i}$ be the noncentral coordinates of
$H_N^ne_0$. The relations
\[
 z_{n+1,i}=v_iA_n+t_iz_{n,i},\qquad
 A_{n+1}=\sum_i v_iz_{n,i},\qquad z_{0,i}=0
\]
give the exact convolution
\begin{equation}\label{conv}
 A_{n+2}=\sum_{k=0}^np_k^{(N)}A_{n-k},\qquad
 p_k^{(N)}=\frac1N\sum_{i=0}^{N-1}(1-t_i)t_i^k.
\end{equation}
These coefficients tend to $1/((k+2)(k+1))$. One may use Riemann sums;
the formal proof instead uses the classical Faulhaber identity~\cite{M}:
\[
 M_N(k):=\frac1N\sum_{i=0}^{N-1}(i/N)^k
 =\sum_{j=0}^k\frac{B_j\binom{k+1}{j}}{k+1}N^{-j}
 \longrightarrow\frac1{k+1}.
\]
Here $B_1=-1/2$. Since $p_k^{(N)}=M_N(k)-M_N(k+1)$, the limit follows.

\newpage
Induction in~\eqref{conv} and~\eqref{rec} proves
$u_n^{(N)}\to u_n$ for every fixed $n$. The strict inequalities
in~\eqref{base} imply that for all sufficiently large $N$,
\[
 b_{12}^{(N)}\le-1/2000,\qquad b_{14}^{(N)}\le b_{12}^{(N)},\qquad
 b_3^{(N)}\le-1/4.
\]
Lemma~\ref{prop} applies with constants independent of $N$.
Passing to the limit gives
\[
 b_{12+2m}\le-\frac1{2000}4^{-m}<0,\qquad b_{2m+1}\le-4^{-m}<0.
\]
Every integer $n\ge11$ is covered. This proves Theorem~\ref{main}.

\section{Formal scope and research status}
The development uses Lean 4.34.1 with mathlib revision
\texttt{d13f23b723b8a846827a245b89c10fc7d3f11612}.
It contains 117 authored declarations in 16 modules. The endpoint
\texttt{FlatBernoulli.source\_conjecture\_319292} has only the hypothesis
$n\ge12$. The finite factorial count, first-block composition bijection,
matrix sums of squares, negative-eigenvalue multiplicity argument, actual
center-moment convolution, coefficient limits, exact rational checks and
strict limiting bounds are all included. The original game-position type
and the source's enumeration theorem are not re-formalized.

The use of finite spectral decomposition, integer-power convexity and
Faulhaber's formula is classical; no new general machinery is claimed.
The contribution is one resolution of the source's explicit all-order conjecture.
The finite checks and supporting lemmas are ingredients of this single result.
The source's ancillary printed equation(23), whose integral sign differs
from the coefficient identity, is not used as a target or counted as an outcome.

The bounded prior-art investigation includes the source, relevant later
papers, reciprocal-log coefficient references, OEIS and public proof
repositories. It cannot certify worldwide priority. Access to the final
JCTA body and one game-theory citer remains limited; no resolution is
inferred from an empty citation corpus. The independent checks and replay instructions accompany the proof files.

\begin{thebibliography}{9}
\bibitem{H} G.~Hetyei, \emph{Decomposition of strongly Bernoulli type games
into elementary kernel factors}, arXiv:0912.0573v1 (2009), especially(16),
Proposition~6.7 and Conjecture~6.10. Published in J.~Combin.~Theory Ser.~A
117(8) (2010), 1107--1126, \href{https://doi.org/10.1016/j.jcta.2010.01.003}{doi:10.1016/j.jcta.2010.01.003}.
\bibitem{M} The mathlib community, \emph{Mathlib}, pinned revision
\texttt{d13f23b7}; \texttt{NumberTheory/Bernoulli.lean},
theorem \texttt{sum\_range\_pow} (classical Faulhaber formula).
\end{thebibliography}
\end{document}
